\documentclass[10pt,twocolumn]{article}

\usepackage[margin=0.72in,columnsep=0.24in]{geometry}
\usepackage[T1]{fontenc}
\usepackage[utf8]{inputenc}
\usepackage{lmodern}
\usepackage{microtype}
\usepackage{amsmath,amssymb}
\usepackage{booktabs}
\usepackage{graphicx}
\usepackage{enumitem}
\usepackage{xcolor}
\usepackage{hyperref}
\usepackage{url}
\usepackage{balance}

\definecolor{linkblue}{HTML}{175A9C}
\hypersetup{
  colorlinks=true,
  linkcolor=linkblue,
  citecolor=linkblue,
  urlcolor=linkblue,
  pdftitle={LogRMS: A Two-Byte Optimizer with TinyStories Training Results},
  pdfauthor={Matthew Belcher}
}
\setlist{nosep,leftmargin=*}
\setlength{\emergencystretch}{1.5em}

\newcommand{\paperauthor}{Matthew Belcher}
\newcommand{\paperaffiliation}{Computer Science, Open Institute of Technology (OPIT), Malta}

\title{\textbf{LogRMS: A Two-Byte Optimizer with\\
TinyStories Training Results}}
\author{\paperauthor\\
\small \paperaffiliation}
\date{13 July 2026}

\begin{document}
\maketitle

\begin{abstract}
Adaptive optimizers often store several values for every model parameter.
LogRMS uses exactly two persistent tensor bytes per parameter: an int8
log-encoded second moment and a uint8 record of recent nonzero gradient signs.
It computes the second-moment recurrence and current update in FP32, then uses
stochastic rounding to store the moment for the next step. Sign history only
reduces updates on coordinates whose gradients alternate frequently. Compared
with ordinary FP32-state AdamW, the persistent tensor state is 75\% smaller.

A controlled 49,515,008-parameter TinyStories systems benchmark processed
471,765,959 training tokens on one NVIDIA L4. Compiled LogRMS achieved a median
full-step time of 107.369 ms versus 107.263 ms for fused AdamW, a 0.10\%
difference, while using 2.0000 versus 8.0000 persistent state bytes per
parameter. Its total wall time was 3.5\% longer, mainly because of a
207-second one-time compilation. Under the shared, untuned schedule, LogRMS
finished at test perplexity 3.7698 versus 3.6911 for fused AdamW, a 2.13\% gap.
Peak live PyTorch allocation was 12.7\% lower, although that number covers the
whole training process rather than the optimizer alone. A separate
5.06M-parameter run and small diagnostic tests show that the same optimizer can
train reliably, but can struggle on curved objectives and rare gradient spikes.
These results support LogRMS as a practical low-state optimizer candidate. They
do not show that it always converges as well as AdamW.
\end{abstract}

\section{Introduction}

Optimizer state can use a large share of training memory.
Ordinary Adam maintains exponential moving averages of both gradients and
squared gradients \cite{kingma2014adam}; in the common FP32-state form these two
tensors require approximately eight bytes per parameter, before weights,
gradients, master copies, or framework metadata. At large model sizes, a few
bytes per parameter add up quickly.

Quantized optimizers reduce this burden by compressing Adam's moments
\cite{dettmers20218bit}, while factorized and structured approaches reduce the
number of stored statistics \cite{shazeer2018adafactor,anil2019sm3}. LogRMS
stores one coordinate-wise RMS statistic in one byte and uses a second byte for
short sign history instead of keeping an Adam-style first moment.

This paper provides:

\begin{enumerate}
  \item an exact two-byte per-parameter persistent tensor layout;
  \item stochastic log encoding that preserves small changes in the stored
        second moment;
  \item a standard \texttt{torch.optim.Optimizer} implementation with normal
        checkpoint support;
  \item a \texttt{torch.compile} execution path; and
  \item transformer results, timing measurements, exact state accounting, and
        diagnostic tests.
\end{enumerate}

The results show that LogRMS can train 5.06M- and 49.5M-parameter transformers.
On the 49.5M model, its compiled implementation had almost the same steady-state
full-step time as fused AdamW. The experiments use one final-run seed at each
scale, so they do not establish general convergence parity, distributed
scaling, or a universal peak-memory advantage.

\section{Related work}

AdaGrad established per-coordinate adaptive scaling from accumulated squared
gradients \cite{duchi2011adaptive}. Adam combines first- and second-moment
exponential averages with bias correction \cite{kingma2014adam}; AdamW separates
weight decay from the loss-gradient update \cite{loshchilov2017decoupled}.
LogRMS keeps Adam's second-moment form and bias correction but removes the first
moment.

Memory-efficient alternatives include Adafactor's factored second moments
\cite{shazeer2018adafactor}, SM3's covering-set accumulators
\cite{anil2019sm3}, and block-wise 8-bit quantization of optimizer state
\cite{dettmers20218bit}. LogRMS keeps a simple coordinate-wise state for tensors
of any rank. Its two bytes are not necessarily the smallest possible
representation when factorization or block metadata is considered.

Sign-based methods demonstrate that gradient direction can remain useful after
discarding magnitude \cite{bernstein2018signsgd}; Lion uses sign information in
an update constructed from momentum-like state \cite{chen2023lion}. LogRMS does
not use a sign-only update. It retains gradient magnitude through RMS
normalization and uses sign history only to reduce selected updates. Stochastic
rounding is commonly used with limited-precision arithmetic
\cite{gupta2015limited}; here it is applied to the stored second moment.

\section{Method}

\subsection{RMS update without a first moment}

For parameter $\theta$, gradient $g_t$, and second-moment decay $\beta_2$, the
optimizer computes
\begin{align}
v_t &= \beta_2 \widetilde v_{t-1} + (1-\beta_2)g_t^2, \\
\widehat v_t &= \frac{v_t}{1-\beta_2^t}, \\
\theta_t' &= (1-\eta\lambda)\theta_{t-1}, \\
\theta_t &= \theta_t' - \eta c_t
\frac{g_t}{\sqrt{\widehat v_t}+\epsilon},
\end{align}
where $\eta$ is the learning rate, $\lambda$ is decoupled weight decay,
$c_t$ is sign confidence, and $\widetilde v_{t-1}$ is the decoded stored state.
All quantities in the recurrence and normalization are calculated in FP32. The
fresh, unquantized $v_t$ is used for the current update; compression affects the
next recurrence.

The state encoding does not limit model precision. LogRMS can update FP32
parameters because it keeps the current gradient and moment calculation in FP32.
Only the stored second moment is compressed for the next step.

\subsection{One-byte logarithmic second moment}

An int8 code $q=-128$ denotes exact zero. Other codes decode as
\begin{equation}
\mathcal D(q) = 2^{q/2}, \qquad q\in\{-127,\ldots,127\}.
\end{equation}
The half-octave spacing covers approximately $2^{-63.5}$ to $2^{63.5}$.
Given adjacent represented values $a\leq v_t\leq b$, encoding selects $b$ with
probability
\begin{equation}
p_b = \frac{v_t-a}{b-a}
\end{equation}
and $a$ otherwise. Consequently, $\mathbb E[\mathcal D(q_t)]=v_t$ away from
range clipping. A stateless integer hash of coordinate, parameter ordinal,
step, and seed supplies the random value without consuming model RNG state or
storing a residual. Stochastic rounding prevents small EMA changes from being
rounded away repeatedly.

\subsection{One-byte sign history}

A uint8 variable-length prefix code stores zero through seven prior nonzero
signs. The current nonzero sign completes an eight-sign sequence. If $T_t$ is
the number of adjacent transitions among those eight signs, the confidence is
\begin{equation}
c_t = 1 - 0.8\frac{T_t}{7}, \qquad T_t\in\{0,\ldots,7\}.
\end{equation}
Thus $c_t\in[0.2,1]$: stable signs receive the ordinary RMS update and frequently
alternating signs are damped. Confidence is exactly one until eight genuine
nonzero observations exist. A zero gradient neither advances history nor counts
as a negative sign.

The implementation uses small cached lookup tensors to decode the history and
map each 8-bit sequence to confidence.

\subsection{Memory accounting}

Table~\ref{tab:payload} reports unique persistent tensor storage held by each
optimizer. Python objects, allocator rounding, and non-tensor metadata are not
included. At one trillion parameters, the difference between two and eight
bytes is about 5.46 TiB in total. Sharding divides that saving among workers.

This number is different from CUDA allocator measurements.
\texttt{max\_memory\_allocated} records the peak live tensor allocation across
the complete PyTorch process, including parameters, gradients, activations,
optimizer state, temporaries, compiler buffers, and evaluation.
\texttt{max\_memory\_reserved} also reflects blocks retained by the
caching allocator. Neither measurement isolates optimizer memory, and both may
omit allocations made outside the PyTorch allocator. We report them as
whole-process measurements, not as optimizer-state measurements.

\begin{table}[t]
\centering
\caption{Persistent tensor state in the 49.5M systems benchmark. Measured
8-bit AdamW includes its tensor metadata in the tested torchao configuration.}
\label{tab:payload}
\footnotesize
\begin{tabular}{@{}lrr@{}}
\toprule
Optimizer & B/parameter & Reduction vs. 8 B \\
\midrule
AdamW, FP32 moments & 8.0000 & --- \\
torchao AdamW 8-bit & 2.0460 & 74.4\% \\
LogRMS log8 & \textbf{2.0000} & \textbf{75.0\%} \\
\bottomrule
\end{tabular}
\end{table}

\begin{table}[t]
\centering
\caption{Aggregate state payload implied by exact 8 and 2 B/parameter layouts.}
\label{tab:scale}
\small
\begin{tabular}{rrrr}
\toprule
Parameters & AdamW & LogRMS & Saving \\
 & (GiB) & (GiB) & (GiB) \\
\midrule
1B & 7.45 & 1.86 & 5.59 \\
7B & 52.15 & 13.04 & 39.12 \\
13B & 96.86 & 24.21 & 72.64 \\
70B & 521.54 & 130.39 & 391.16 \\
1T & 7,450.58 & 1,862.65 & 5,587.94 \\
\bottomrule
\end{tabular}
\end{table}

\section{5.06M TinyStories convergence study}

\subsection{Protocol}

We trained a 5,063,424-parameter causal transformer on the official TinyStories
dataset \cite{eldan2023tinystories}. Each final optimizer run saw exactly two
complete training epochs: 943,531,918 tokens in 115,
